\section{Priorités de calcul}

\begin{rappelbox}
Dans un calcul, on effectue les opérations dans cet ordre :

\begin{enumerate}
    \item les calculs entre parenthèses ;
    \item les multiplications et divisions ;
    \item les additions et soustractions.
\end{enumerate}
\end{rappelbox}

\begin{exemplebox}
\[
3+(-4)\times(-2)
\]

On commence par la multiplication :

\[
3+8=\answer{11}
\]
\end{exemplebox}

\begin{exemplebox}
\[
(-5+2)\times(-4)
\]

On commence par la parenthèse :

\[
(-3)\times(-4)=\answer{12}
\]
\end{exemplebox}

\begin{attentionbox}
On ne calcule pas simplement de gauche à droite lorsque le calcul
contient plusieurs opérations.

\[
3+2\times5 \neq (3+2)\times5
\]
\end{attentionbox}
